\RequirePackage{pdfmanagement}
\documentclass[a4paper]{l3doc}

\usepackage{microtype}
\usepackage[american, bidi=basic]{babel}
% \babelprovide[import, onchar=ids]{polytonicgreek}
% \babelprovide[import, onchar=ids]{russian}
% \babelprovide[import, onchar=ids fonts]{armenian}
% \babelprovide[import, onchar=ids fonts]{hebrew}
% \babelprovide[import, onchar=ids fonts]{syriac}
\babelfont{rm}[Numbers=Lining]{Brill}
\babelfont{tt}[Scale=MatchLowercase]{DejaVu Sans Mono}
% \babelfont[armenian]{rm}[Scale=0.85]{Noto Serif Armenian}
% \babelfont[hebrew]{rm}[Scale=MatchLowercase, Contextuals=Alternate]{SBL Hebrew}
% \babelfont[syriac]{rm}[Scale=MatchLowercase]{Estrangelo Edessa}
% \babelfont[hebrew]{tt}[Scale=MatchLowercase]{Miriam Libre Medium}
% \babelfont[syriac]{tt}[Scale=0.85]{Noto Sans Syriac}
\usepackage{csquotes}
\usepackage[style=sbl]{biblatex}
\addbibresource{biblatex-sbl.bib}

% allow lining figures in URLs
\renewcommand*{\url}[1]{\href{#1}{#1}}

% SBL style footnotes
\usepackage[bottom, ragged]{footmisc}
\makeatletter
\renewcommand*{\@makefntext}[1]{%
  \parindent \footnotemargin
  \@thefnmark . \enspace
  \footnotelayout #1}
% redefine \maketitle to better match \thanks footnote
\renewcommand\maketitle{\par
  \begingroup
    \renewcommand\thefootnote{\@fnsymbol\c@footnote}%
    \renewcommand*{\@makefntext}[1]{%
      \parindent \footnotemargin
      \@textsuperscript{\@thefnmark}%
      \footnotelayout ##1}%
    \if@twocolumn
      \ifnum \col@number=\@ne
        \@maketitle
      \else
        \twocolumn[\@maketitle]%
      \fi
    \else
    \newpage
      \global\@topnum\z@   % Prevents figures from going at top of page.
      \@maketitle
    \fi
    \thispagestyle{plain}\@thanks
  \endgroup
  \setcounter{footnote}{0}%
  \global\let\thanks\relax
  \global\let\maketitle\relax
  \global\let\@maketitle\relax
  \global\let\@thanks\@empty
  \global\let\@author\@empty
  \global\let\@date\@empty
  \global\let\@title\@empty
  \global\let\title\relax
  \global\let\author\relax
  \global\let\date\relax
  \global\let\and\relax
}
\makeatother

\begin{document}

\title{The \pkg{biblatex-sbl} Package}
\author{David Purton\thanks{Email: \url{dcpurton@marshwiggle.net}}}
\date{2026-03-17 v2.0}

\maketitle

\begin{abstract}
  This package provides \pkg{biblatex} support for citations, bibliography and
  abbreviations in the format specified by the second edition of the
  \emph{Society of Biblical Literature (SBL) Handbook of Style} including
  updates on the \emph{SBL Handbook of Style} blog.
\end{abstract}

\tableofcontents

\section{Introduction}

\pkg{biblatex-sbl} provides support to \pkg{biblatex} and LaTeX for citations,
bibliography and abbreviations in the style recommended by the Society of
Biblical Literature (SBL). Shorthand citations and a list of abbreviations
containing journals, series, shorthands and ancient works are handled
automatically. Only note style citations, not author-date citations, are
supported. \pkg{biblatex-sbl} is compatible with \pkg{biblatex}'s support for
\pkg{hyperref}.

The style seeks to support the latest requirements outlined on the \emph{SBL
Handbook of Style} blog \autocite{sblhsblog}. All examples from the
handbook, student supplement, and blog are fully supported (noting that the
blog updates many recommendations).

There is also sophisticated handling of subsequent citations, reprints, and
the many unique entry types required by SBL style, such as lexicons,
dictionaries, and encyclopaedias.

\pkg{biblatex-sbl} is compatible with \pkg{biblatex}'s \pkg{hyperref} and
\pkg{babel} support, although custom strings are only provided for English at
this stage.

A companion bundle of packages and classes are distributed separately to
\pkg{biblatex-sbl}. These provide support for referencing Bible passages
(\pkg{bibleref-sbl}), font set up (\pkg{sblfonts}), indexing (\pkg{sblidx}),
term papers (\pkg{sblarticle}), theses (\pkg{sblreport}) and print ready books
(\pkg{sblbook}).

This manual only covers custom features of \pkg{biblatex-sbl}. For general
help with \pkg{biblatex} you should consult its extensive manual.

\subsection*{Note}

The blog has now made SBL style a rolling style that changes with each blog
post. Because of this, preserving backward compatibility between releases is
not feasible for time I have available to maintain this style and neither is
it possible to provide options to vary the style to cater for different
seminaries and theological colleges.

\textbf{In particular, version 2.0 and later of \pkg{biblatex-sbl} completely
breaks compatibility with the version 0.x releases.} If you need to access a
legacy release, you should download it from the GitHub repository and place
the style files locally in your project.\footnote{See
\url{https://github.com/dcpurton/biblatex-sbl/releases} for legacy releases
and \url{https://github.com/dcpurton/biblatex-sbl/tree/legacy} for legacy
source code.} Legacy 0.x versions are no longer maintained or supported.

\end{document}
